% SixtySecondsCV LaTeX template
% Copyright © 2020 L. Gueroult <lgueroult-at-hotmail.com>
% Licensed under the 3-Clause BSD License. See LICENSE file for details.
%
% Please visit https://github.com/LaGuer/SixtySecondsCV for the most
% recent version! For bugs or feature requests, please open a new issue on
% github.
%
% Contributors
% ------------
% * ifokkema
% * Bertbk
% * Hespe
% * esben
%
% Attributions
% ------------
% * sixtysecondscv is based on the fortysecondcv class by Rene Wirnata
%   (rene.wirnata@pandascience.net), released under the 3-Clause BSD license and available at
%   https://github.com/PandaScience/FortySecondsCV/
% * fortysecondscv is based on the twentysecondcv class by Carmine Spagnuolo
%   (cspagnuolo@unisa.it), released under the MIT license and available at
%   https://github.com/spagnuolocarmine/TwentySecondsCurriculumVitae-LaTex
% * further attributions are indicated immediately before corresponding code


%-------------------------------------------------------------------------------
%                             ADDITIONAL PACKAGES
%-------------------------------------------------------------------------------
\documentclass[
	a4paper,
	% showframes,
	% vline=2.2em,
	%maincolor=cvgreen,
	sidecolor=gray!8,
	sectioncolor=materialblue,
	subsectioncolor=materialindigo,
	% itemtextcolor=black!80,
	sidebarwidth=0.27\paperwidth,
	% topbottommargin=0.03\paperheight,
	% leftrightmargin=20pt,
	profilepicsize=3.5cm,
	% profilepicborderwidth=3.5pt,
	profilepicstyle=profilecircle,
	profilepiczoom=1.15,
	% profilepicxshift=0mm,
	% profilepicyshift=0mm,
	% profilepicrounding=1.0cm,
	% logowidth=4.5cm,
	% logospace=5mm,
	% logoposition=before,
]{sixtysecondscv}

%-------------------------------------------------------------------------------
%                             ADDITIONAL PACKAGES
%-------------------------------------------------------------------------------
% improve word spacing and hyphenation
\usepackage{microtype}
\usepackage{ragged2e}

% uncomment in case you don't want any hyphenation
% \usepackage[none]{hyphenat}

% take care of proper font encoding
% \ifxetexorluatex
% 	\usepackage{fontspec}
% 	\defaultfontfeatures{Ligatures=TeX}
%	\newfontfamily\headingfont[Path = fonts/]{segoeuib.ttf} % local font
% \else
	\usepackage[utf8]{inputenc}
	\usepackage[T1]{fontenc}
%	\usepackage[sfdefault]{noto} % use noto google font
% \fi

% enable mathematical syntax for some symbols like \varnothing
\usepackage{amssymb}

% bubble diagram configuration
\usepackage{smartdiagram}
\smartdiagramset{
	% default font size is \large, so adjust to harmonize with sidebar layout
	bubble center node font = \footnotesize,
	bubble node font = \footnotesize,
	% default: 4cm/2.5cm; make minimum diameter relative to sidebar size
	bubble center node size = 0.4\sidebartextwidth,
	bubble node size = 0.25\sidebartextwidth,
	distance center/other bubbles = 1.5em,
	% set center bubble color
	bubble center node color = maincolor!70,
	% define the list of colors usable in the diagram
	set color list = {materiallime, materialamber,
	materialcyan, materialblue, materialindigo},
	% sets the opacity at which the bubbles are shown
	bubble fill opacity = 0.8,
}

%-------------------------------------------------------------------------------
%                            PERSONAL INFORMATION
%-------------------------------------------------------------------------------
%% mandatory information
% your name
\cvname{Antonio Rizza}
% job title/career
\cvjobtitle{Postdoctoral \\ \\Researcher \\ \\
% \small Date of birth: 18/12/1995\\
% Place of birth: Padova (PD)\\
% Tax code: RZZNTN95T18G224W
}

\write18{wget https://raw.githubusercontent.com/gosquared/flags/master/flags/flags/shiny/64/Germany.png}
\write18{wget https://raw.githubusercontent.com/gosquared/flags/master/flags/flags/shiny/64/France.png}
\write18{wget https://raw.githubusercontent.com/gosquared/flags/master/flags/flags/shiny/64/Spain.png}
\write18{wget https://raw.githubusercontent.com/gosquared/flags/master/flags/flags/shiny/64/United-States.png}
\write18{wget https://raw.githubusercontent.com/gosquared/flags/master/flags/flags/shiny/64/Russia.png}
\write18{wget https://raw.githubusercontent.com/gosquared/flags/master/flags/flags/shiny/64/Switzerland.png}
\write18{wget https://ipfs.io/ipfs/QmZiD898vX1yaukukNbgbuSoUtqVRD262k4qK4xWB523LM?filename=QmZiD898vX1yaukukNbgbuSoUtqVRD262k4qK4xWB523LM.png}


%% optional information
% profile picture
\cvprofilepic{antonio_rizza_cut.jpg}
% logo picture
%\cvlogopic{pics/logo_txt.png}

% NOTE: ordering in sidebar will mimic the following order
% date of birth
% \cvbirthday{\today}
% short address/location, use \newline if more than 1 line is required

% ADD ADDRESS HERE

% \cvaddress{Viale G. Grassi 113, Lecce, Italy}
% phone number
% \cvphone{(+39) 3209218918}
% skype ID
% \cvskype{myskypeID}
% personal website
% \cvsite{https://laguer.github.io/sixtysecondscv/}
% email address

\cvmail{arizza@stanford.edu}


% \cvmail{antonio.rizza.mail@gmail.com}
% pgp key
% \cvkey{4096R/FF00FF00}{0xAABBCCDDFF00FF00}
% any other custom entry
% \cvcustomdata{}}
\cvsite{https://antoniorizza.github.io}{Personal Website}
\cvlinkedin{antoniorizza95}
\cvscholar{https://scholar.google.com/citations?user=JqvWV-AAAAAJ&hl=it}{Google Scholar}

%-------------------------------------------------------------------------------
%                              SIDEBAR 1st PAGE
%-------------------------------------------------------------------------------
% add more profile sections to sidebar on first page
\addtofrontsidebar{
	% include gosquare national flags from https://github.com/gosquared/flags;
	% naming according to ISO 3166-1 alpha-2 country codes
	\graphicspath{{pics/flags/shiny/}}
	\vspace{2ex}

% 	% social network accounts incl. proper hyperlinks
% 	\profilesection{Social Network}
% 		\begin{icontable}{2.5em}{1em}
% 			\social{\aiOverleafSquare}
% 				{https://www.overleaf.com/latex/templates/SixtySecondsCV/}
% 				{Overleaf Link}
% 			\social{\aiResearchGate}
% 				{https://www.researchgate.net/profile/}
% 				{Researchgate Link}
% 			\social{\faGithub}
% 				{https://github.com/LaGuer/SixtySecondsCV}
% 				{Github Project Page Link}
% 		\end{icontable}

	\profilesection{Languages}
	\normalsize
	Italian, English, Spanish, French \hspace{3cm}
	\vspace{0.01cm}
        \profilesection{Expertise}

    \vspace{0.01cm}
   Spacecraft GNC\\

   Astrodynamics and Formation\\ Flying\\ 
      
   Modeling and Simulation of Space Systems \\

   Dynamical Systems Theory\\

   Optimal Control Theory\\
   
   Non-linear State Estimation\\

	\vspace{0.01cm}
    \profilesection{Computer Skills}
    % \large \textbf{Confident}\\ \normalsize
    \vspace{0.01cm}
	MATLAB/Simulink\\

    C++, Python (VSCode)\\
    
    NASA tools (SPICE, GMAT)\\
    
    CAD tools (SolidWorks, Solid Edge, Inventor)\\
    
    % Editors/Styles (LaTeX, Office, VSCode)\\

    % OS  (iOS, Linux, Windows)
	\vspace{0.01cm}


		{\footnotesize \color{black!70} \hfill 

% 	In compliance with the European Regulation 679/2016, I hereby authorize you to use and process my personal data
}
    %\pointskill{\flag{Russia.png}}{Russia}{5}
    
%     \begin{figure}\centering
% 		\smartdiagram[bubble diagram]{
% 			\textcolor{white}{\textbf{Being a}} \\
% 			\textcolor{white}{\textbf{Panda}}, % center bubble
% 			\textcolor{black!90}{Eating},
% 			\textcolor{black!90}{Sleeping},
% 			\textcolor{black!90}{Rolling},
% 			\textcolor{black!90}{Playing},
% 			\textcolor{black!90}{Chilling}
% 		}
% 	\end{figure}
	
% 	\profilesection{Barskills}
% 	\barskill{\faSkyatlas}{Wearing asian rice hats}{60}
% 	\barskill{\faImage}{Playing Chess}{30}
% 	\barskill{\faMusic}{Playing the bamboo stick}{50}
}
    

%-------------------------------------------------------------------------------
%                              SIDEBAR 2nd PAGE
%-------------------------------------------------------------------------------
\addtobacksidebar{

}


%-------------------------------------------------------------------------------
%                         TABLE ENTRIES RIGHT COLUMN
%-------------------------------------------------------------------------------
\begin{document}

\makefrontsidebar


\cvsection{Headline}
% \begin{center}
\normalsize

I am a Spacecraft Guidance Navigation and Control (GNC) Engineer with experience in modeling, simulation and testing of dynamical Space Systems, spacecraft autonomy, precise formation flying and astrodynamics. I had leading GNC roles for Low Earth Orbit (LEO) and deep-space missions sponsored by national agencies. At Stanford, I am currently focused on enabling GNC solutions for the next generation of distributed space systems. I have been GNC Engineer for the \href{https://tyvak.eu/missions/milani/}{\textcolor{materialblue}{Milani}}, \href{https://www.esa.int/Enabling_Support/Space_Engineering_Technology/Shaping_the_Future/LUMIO_New_CubeSat_Illuminating_Lunar_Impacts}{\textcolor{materialblue}{LUMIO}}, and \href{https://starimission.space/index.html}{\textcolor{materialblue}{STARI}} mission.

%   Enthusiastic and motivated Space Engineer, with strong background in GNC, System Engineering and ground testing methodologies for Aerothermal Demise. Currently enrolled in a PhD course at Politecnico di Milano.

% \end{center}
\cvsection{Research Experience}
\begin{cvtable}[1.5]
\cvitem{\mbox{01/2025 - ongoing}}{Postdoctoral Scholar}{Stanford University - SLAB team}{Enabling safe, robust and high precision GNC solutions for the next generation of distributed space systems. Lead GNC engineer for the \href{https://starimission.space/index.html}{\textcolor{materialblue}{STARI}} mission.}

\cvitem{\mbox{11/2024 - 12/2024}}{Research Fellow (Senior)}{Politecnico di Milano - DART team}{Developing guidance and control solution for proximity opeartions around small bodies. Supporting the development of a robotic facility for validation and verification of GNC algorithm.}

\cvitem{\mbox{01/2023 - 12/2023}}{Visiting Research Scholar}{Stanford University - SLAB team}{Goal-oriented trajectory optimization about asteroids using sequential convex programming. Hardware in the loop validation of Active SLAM approaches using the robotic testbed for rendezvous and optical navigation at SLAB.}

\cvitem{\mbox{11/2021 - 11/2024}}{PhD (\textit{cum laude}), Aerospace Eng.}{Politecnico di Milano - DART team}{Goal-oriented guidance under uncertainties for small bodies exploration. GNC engineer for the \href{https://tyvak.eu/missions/milani/}{\textcolor{materialblue}{Milani}} and \href{https://www.esa.int/Enabling_Support/Space_Engineering_Technology/Shaping_the_Future/LUMIO_New_CubeSat_Illuminating_Lunar_Impacts}{\textcolor{materialblue}{LUMIO}} missions.}

% guidance methodologies for on-board real-time implementation in deep-space CPOD scenarios.}
\cvitem{\mbox{11/2020 - 10/2021}}{Research Fellow}{\normalsize \textit{Politecnico di Milano - DART team}}{Development of high-fidelity six-degrees-of-freedom simulator for autonomous deep-space missions.}

% \cvitem{10/2019 - 06/2020}{M.Sc. Thesis}{\normalsize \textit{VKI, Belgium} }{Simulation and testing of gas-surface interaction phenomena in \\ atmospheric re-entry conditions for Design for Demise analysis.}

% \cvitem{08/2017 - 09/2017}{Research Internship}{\normalsize \textit{VKI, Belgium}}{
% 		Software developer for the VKI Expandable Stability and Transition Analysis toolkit to predict boundary layer transition in high speed flow.}
\end{cvtable}
% \vspace{0.5cm}

\cvsection{Education}
% \cvsubsection{Postgraduate Training}
\begin{cvtable}[1.5]
\cvitem{\mbox{09/2017 - 06/2020}}{M.Sc. \textit{cum laude} in Space Engineering }{\normalsize \textit{Politecnico di Milano}}{}
% \cvitem{05/2019}{\textbf{ESA Accademy Space Debris Training Course}}{\normalsize \textit{ESEC GALAXIA, Belgium}}{}
\cvitem{\mbox{10/2014 - 09/2017}}{Bachelor's in Aerospace Engineering}{\normalsize \textit{Politecnico di Milano}}{}
\end{cvtable}
% \vspace{0.5cm}

% \begin{cvtable}[1.5]
% 	\cvitem{}{Spacecraft Attitude Dynamics and Control Course}{\normalsize \textit{Politecnico di Milano}}{}
% 	\cvitem{}{Spacecraft Attitude Dynamics and Control Course}{\normalsize \textit{Politecnico di Milano}}{}
% % \cvitem{09/2021-12/2021}{Spacecraft Attitude Dynamics and Control Course}{Politecnico di Milano}
% \end{cvtable}

\cvsection{Research projects and Missions} 
% \cvsubsection{Postgraduate Training}
\begin{cvtable}[1.5]
\cvitem{01/2026-ongoing}{Flight Demonstration of Autonomous Navigation in GEO}{}{Collaboration with Blue Origin to demonstrate in  GEO for the first time autonomous navigation capabilities based on angles only navigation and CNN-based relative pose estimation to known targets.}

\cvitem{01/2025-ongoing}{FALCON - Research Co-lead}{}{Fast Autonomous Lost-in-space Catalog-based Optical Navigation (\href{https://slab.stanford.edu/projects/falcon}{\textcolor{materialblue}{FALCON}}) is a new alternative space-based Positioning, Navigation and Timing (PNT) system for responsive Space Situational Awareness (SSA) of on-orbit assets.}

\cvitem{01/2025-ongoing}{STARI - Orbit and GNC Engineer - Phase A to E}{}{The STarlight Acquisition and Reflection toward Interferometry (\href{https://starimission.space/index.html}{\textcolor{materialblue}{STARI}}) mission is a technology demonstration designed to advance key technologies for future space-based interferometry. The mission makes use of two 8U CubeSats flying in close formation, effectively simulating one arm of a space interferometer.}

\cvitem{09/2022-12/2024}{LUMIO - AOCS Engineer - Phase B}{}{The Lunar Meteoroid Impact Observer (\href{https://dart.polimi.it/projects/}{\textcolor{materialblue}{LUMIO}}) is a CubeSat mission to a quasi-halo orbit at the Earth-Moon L2 Lagrangian point designed to complement ground based observations and characterize the lunar meteoroid environment.}

\cvitem{11/2020-06/2022}{Milani - GNC Engineer - Phase A to D}{}{The (\href{https://dart.polimi.it/projects/}{\textcolor{materialblue}{Milani}}) CubeSat is part of ESA's Hera mission, the European contribution to the ESA-NASA collaboration AIDA (Asteroid Impact \& Deflection Assessment) around the binary asteroid 65803 Didymos.}
\end{cvtable}
% \vspace{0.5cm}

% \cvsection{Teaching and tutoring}

% Grading and tutoring graduate students for the course AA279D \textit{Control of Distributed Space Systems} at Stanford University.\\

% Teaching assistant for the graduate course in \textit{Spacecraft Attitude Dynamics and Control} at Politecnico di Milano for two years.\\

% Co-advisor of PhD, master and undergraduate students at the Space Rendezvous Laboratory (SLAB) at Stanford University. \\

% Co-advisor of tens of M.Sc. students during their final thesis work at the Deep-space Astrodynamics Research and Technology (DART) group at Politecnico di Milano.\\

% \cvsection{Publication records}

% \url{https://scholar.google.com/citations?user=JqvWV-AAAAAJ&hl=it}
% \vspace{0.5cm}

% \cvsection{Journals}

% \textbf{A. Rizza}, C. Giordano, F. Topputo. "Goal-Oriented Guidance Strategy for Binary Asteroid Exploration".  Astrodynamics (In-press-2024) .\\

% \textbf{A. Rizza}, S. D'Amico, F. Topputo. "Goal-oriented Guidance under Uncertainties for ASteroid Mapping using Sequential Convex Programming". Under review for Journal of Guidance Control \& Dynamics (JGCD) (2024).\\

% M. Pugliatti, F. Piccolo, \textbf{A. Rizza}, V. Franzese, F. Topputo. "The vision-based guidance, navigation, and control system of Hera's Milani Cubesat". Acta Astronautica (2023).\\

% % \vspace{0.5cm}

% \cvsection{Conferences}

% \textbf{A. Rizza}, S. D'Amico, F. Topputo, “Goal-oriented guidance under uncertainties for small bodies exploration". 2025 AAS/AIAA Space Flight Mechanics Meeting,
% Kaua'i, Hawaii, January 19-23 2025 (SFMM25). \\


% \textbf{A. Rizza}, C. Bonagura, P. Panicucci, F. Topputo, “Modular Pipeline for Small Bodies Gravity Field Modeling: Enhancing accuracy and efficiency for proximity operations". Accepted for presentation at the 75th International Astronautical Congress 2024 (IAC2024). \\

% \textbf{A. Rizza}, F. Topputo, “Gradient based reachability analysis for goal-oriented guidance. Application to Eros proximity operations,” 6th International Workshop on Key Topics in Orbit Propagation Applied to SSA, Arras, June 12-14, 2024 (KePassa2024). \\

% \textbf{A. Rizza}, F. Topputo, and S. D’Amico, “Goal-oriented asteroid mapping under uncertainties using Sequential Convex Programming,” AIAA SCITECH 2024 Forum, (SciTech2024). \\

% \textbf{A. Rizza}, C. Giordano, F. Topputo. "Goal-Oriented Guidance Strategy for Binary Asteroid Exploration: Investigation of Different Metrics." In 45th AAS Guidance, Navigation and Control Conference (GN\&C 2023).\\

% \textbf{A. Rizza}, F. Piccolo, P. Panicucci, S. Borgia, G. Saita, G. Fumo, G. Merisio, L. Provinciali, F. Topputo. "Design, Analysis and Validation of the ADCS for the LUMIO mission." 74th International Astronautical Congress 2023 (IAC2023).\\

% P. Panicucci, F. Piccolo, S. Borgia, \textbf{A. Rizza}, V. Franzese, & F. Topputo. (2023). "Current Status of the LUMIO Autonomous Optical Navigation Experiment". In the 12th International Conference on Guidance, Navigation \& Control Systems (GNC) and 9th International Conference on Astrodynamics Tools and Techniques (ICATT).\\

% \textbf{A. Rizza}, F. Piccolo, M. Pugliatti, P. Panicucci, F. Topputo, "Hardware-in-the-loop simulation framework for CubeSats proximity operations: application to the Milani mission". 73rd International Astronautical Congress (IAC), Paris, France, 18-22 September 2022. (IAC2022)\\

% M. Pugliatti, \textbf{A. Rizza}, F. Piccolo, V. Franzese, C. Bottiglieri, C. Giordano, F. Ferrari, F. Topputo. "The Milani mission: overview and architecture of the optical-based GNC system". In 2022 AAS San Diego (AAS 2022).\\

% M. Pugliatti, V. Franzese, \textbf{A. Rizza}, F. Piccolo, C. Bottiglieri, C. Giordano, F. Ferrari, F. Topputo. "Design of the on-board image processing of the milani mission". In 2022 AAS/AIAA Astrodynamics Specialist Conference (AAS 2022).\\

% C. Bottiglieri, F. Piccolo, \textbf{A. Rizza}, C. Giordano, M. Pugliatti, V. Franzese, F. Ferrari, F. Topputo. "Trajectory design and orbit determination of Hera's Milani CubeSat". In 2021 AAS/AIAA Astrodynamics Specialist Conference (AAS 2021).\\

% \textbf{A. Rizza}, M. Pugliatti, F. Piccolo, V. Franzese, C. Bottiglieri, C. Giordano, F. Ferrari, F. Topputo. "A semi-autonomous optical-based GNC design for Milani mission". In 12th European CubeSat Symposium (ECS 2021).

% \cvsection{Academic projects}
% % \cvsubsection{Postgraduate Training}
% \begin{cvtable}[1.5]
% %	\cvitem{08/2019}{Mathematical Programming for Space Debris orbit determination}{\normalsize \textit{}}
% %	{\\Optimization of the observation planning for a ground station with the aim of minimizing the error on debris orbit determination.}	
% 	\cvitem{03/2019 - 07/2019}{Lunar tech kit, a preliminary design (Team)}{\normalsize \textit{}}
% 	{Played the role of Mission Analysis Engineer in a team project for the phase A, preliminary mission design of a lunar science tech kit mission to map lunar resources and distribution of lava tubes.}
% 	\cvitem{10/2018 - 01/2019}{\textbf{Spacecraft Attitude Dynamics}}{\normalsize \textit{}}
% 		{Design and simulation of the ADCS for a ComSat in GEO.}
% 	\cvitem{12/2018 - 01/2019}{Model based design on GOCE (Team)}{\normalsize \textit{}}
% 	{Modeling and Simulation of the Drag Free and Attitude Control System for GOCE Spacecraft. Multi-objective optimization of the control parameters and out of nominal robustness verification.}
% % 	\cvitem{10/2018}{ReDSHIFT Space Debris Day}{\normalsize \textit{}}
% % 	{One day workshop on the European project ReDSHIFT on de-orbiting and disposal of space debris.}	
% 	\cvitem{10/2017 - 01/2018}{Mission to Uranus (Team)}{\normalsize \textit{}}
% 	{Design of optimal interplanetary trajectory from Earth to Uranus with gravity assist on Mars. Robustness verification against an high fidelity 10 bodies simulation. }
% % 	Mixed grid - heuristic optimization of an interplanetary trajectory from Earth to Uranus performing a FlyBy around Mars. Comparison of the patched conics approach against a 10 body simulation of the Solar System.
% % 	}
% 	\cvitem{10/2017 - 01/2018}{Perturbed Earth environment (Team)}{\normalsize \textit{}}
% 	{Long term orbit propagation in MEO under the effects of orbital perturbations.}	
% % 	\cvitem{09/2016 - 01/2017}{A Space Propulsion trade off}{\normalsize \textit{}}
% % 	{Preliminary analysis and design of a suitable propulsion system to perform an Hohmann transfer in Earth orbit, comparison between storable liquid and electric solutions.}	
% \end{cvtable}

% \vspace{1cm}


% All the certificates and documentation related to language knowledge and experiences abroad can be made available on request.

%Autorizzo al trattamento dati ai sensi del GDPR 2016/679 del 27 aprile 2016 (Regolamento Europeo
%relativo alla protezione delle persone fisiche per quanto riguarda il trattamento dei dati personali).
%Autorizzo la pubblicazione del Curriculum Vitae sul sito istituzionale del Politecnico di Milano (sez.
%Amministrazione Trasparente) in ottemperanza al D. Lgs n. 33 del 14 marzo 2013 (e s.m.i.).

% \newpage
% \makebacksidebar



% \cvsubsection{Study}
% \begin{cvtable}[1.5]
% 	\cvitem{2006 -- 2008}{Master Studies Panda Science}{Panda Academy}
% 		{Focus: Advanced rice hat studies and nouveau rain-reflecting cover
% 		materials.}
% 	\cvitem{}{Master Theses ($\varnothing\, 1,0$)}{CRSW Institute}
% 		{Impact of solar radiation onto rice hat cover materials with special
% 		attention to water resistance.}
% 	\cvitem{2003 -- 2006}{Bachelor Studies PandaScience}{Panda Academy}
% 		{Focus: Bamboo stick morphology and its usage in different instruments.}
% 	\cvitem{}{Bachelor Theses ($\varnothing\, 1,0$)}{BambooStick  Institute}
% 		{The bamboo stick: An underestimated instrument in orchestras?}
% \end{cvtable}

% \cvsection{Publications}
% \begin{cvtable}
% 	\cvpubitem{Back to Cosmos}{Sanchez et Al.}
% 		{Progress in Physics}{2019}
%     \cvpubitem{Science unification and Permanent Cosmology}{Sanchez et Al.}
%         {Journal de Mathematiques pures et appliquees}{2020}
% \end{cvtable}

% \cvsection{Awards}
% \begin{cvtable}
% 	\cvitem{2010 -- now}{Tiger of the Year}{OBS}{}
% 	\cvitem{2001 -- now}{Voila Engine Search}{H3D}{}
% \end{cvtable}


% \cvsection{Extra-Curricular Activities}
% \begin{cvtable}
% 	\cvitemshort{Sports}{Sailing, Golf}
% 	\cvitemshort{Music}{Drums}
% 	\cvitemshort{Education}{Machine Learning, CNN, AI}
% \end{cvtable}


% \newpage
% \makebacksidebar
% % \newgeometry{
% % 	top=\topbottommargin,
% % 	bottom=\topbottommargin,
% % 	right=\leftrightmargin,
% % 	left=\leftrightmargin
% % }

% \cvsection{section}
% \cvsubsection{Subsection}
% \begin{cvtable}
% 	\cvitem{<dates>}{<cv-item title>}{<location>}{<optional: description>}
% \end{cvtable}

% \cvsection{cvitem}
% \cvsubsection{Multi-line with longer description}
% \begin{cvtable}
% 	\cvitem{date}{Description}{location}{Some longer and more detailed
% 		description, that takes two lines of space instead of only one.}
% 	\cvitem{date}{Description}{location}{Some longer and more detailed
% 		description, that takes two lines of space instead of only one.}
% 	\cvitem{date}{Description}{location}{Some longer and more detailed
% 		description, that takes two lines of space instead of only one.}
% \end{cvtable}

% \cvsubsection{One-line without description}
% \begin{cvtable}
% 	\cvitem{Award}{One-line description}{Sponsor}{}
% 	\cvitem{Award}{One-line description}{Sponsor}{}
% 	\cvitem{Award}{One-line description}{Sponsor}{}
% \end{cvtable}

% \cvsection{cvitemshort}
% \cvsubsection{One-line}
% \begin{cvtable}
% 	\cvitemshort{Key}{Some further description}
% 	\cvitemshort{Key}{Some further description}
% 	\cvitemshort{Key}{Some further description}
% \end{cvtable}

% \cvsubsection{Multi-line with longer description}
% \begin{cvtable}
% 	\cvitemshort{Key}{Some further description. Can fill even more than
% 		only one single line while still keeping the correct indendation level.}
% 	\cvitemshort{Key}{Some further description. Can fill even more than
% 		only one single line while still keeping the correct indendation level.}
% 	\cvitemshort{Key}{Some further description. Can fill even more than
% 		only one single line while still keeping the correct indendation level.}
% \end{cvtable}

% \cvsignature 
%  \watermark.
%\begin{Form}
%  \digsigfield{5cm}{3cm}{\cvsignature}
%\end{Form}



\end{document}
